\subsection{严格态射与真性}

引理 1. --- 设 E 和 F 是拓扑向量空间，u 是从 E 到 F 中的连续线性映射，U 是 E 中 0 的一个邻域。假设 u 诱导一个从 U 到 F 的闭子集上的同胚。则 u 的像 I 是闭的，且 u 诱导一个从 E 到 I 上的同胚。

集 \(u(\mathrm{U})\) 是 I 中 0 的一个邻域；它在 F 中是闭的，故 F 的子群 I 在 0 处局部闭，从而它是闭的（TG, III, p. 7, 命题 4）。

由于 \(\operatorname{Ker}(u)\) 与 U 只交于 0 这一点，有 \(\operatorname{Ker}(u) = \{0\}\) ，故映射 \(u\) 是单射。设 V 是包含于 U 的 E 中 0 的闭邻域。由于 \(u(\mathring{\mathrm{V}})\) 是 \(u(\mathrm{V})\) 中 0 的一个邻域，存在 F 中 0 的均衡邻域 W，使得 \(u(\mathrm{V}) \cap \mathrm{W} \subset u(\mathring{\mathrm{V}})\) 。集 \(u^{-1}(\mathrm{W})\) 是均衡的，故是连通的。它包含 0 且包含于 \(\mathring{\mathrm{V}} \cup (\mathrm{E} - \mathrm{V})\) 。由于 \(\mathring{\mathrm{V}}\) 与 E-V 是 E 的不相交开子集，集 \(u^{-1}(\mathrm{W})\) 包含于 \(\mathring{\mathrm{V}}\) ，从而 \(\mathrm{W} \cap \mathrm{I} \subset u(\mathring{\mathrm{V}})\) 。于是 \(u(\mathring{\mathrm{V}})\) 是 I 中 0 的一个邻域。这蕴含 \(u\) 诱导一个从 E 到 I 上的同胚。

命题 1. --- 设 E 是分离局部凸空间，F 是局部凸空间，u 是从 E 到 F 中的连续线性映射。下列条件等价：

\begin{enumerate}
\def\labelenumi{(\roman{enumi})}
\item
  映射 u 是严格态射，它的核有限维且它的像在 F 中是闭的；
\item
  存在 E 中 0 的闭邻域 V，使得 u 在 V 上的限制是真映射（TG, I, p. 72）。
\item
  \(\Longrightarrow\) (ii)：假设 u 满足条件 (i)。由于 u 的核有限维，它在 E 中有一个拓扑补 \(E_{1}\) （III, p. 55, 命题 1），且在 \(\operatorname{Ker}(u)\) 中存在 0 的一个紧邻域 C。把 E 与 \(E_{1} \times \operatorname{Ker}(u)\) 等同起来；此时集 \(V = E_{1} \times C\) 是 E 中 0 的一个闭邻域。u 在 V 上的限制是 \(E_{1} \times C\) 到 \(E_{1}\) 上的投影（它是真映射，TG, I, p. 77, 推论 5）与 \(u_{1}\) （即 u 在 \(E_{1}\) 上的限制）的复合。而 \(u_{1}\) 是从 \(E_{1}\) 到 F 的闭子空间上的一个同胚，故是真映射（TG, I, p. 72, 命题 2）。两个真映射的复合是真映射（TG, I, p. 73, 命题 5, a)），所以 u 在 V 上的限制是真映射。
\item
  \(\Longrightarrow\) (i)：设 V 是 E 中 0 的一个闭邻域，使得 u 在 V 上的限制 v 是真映射。此时集 \(\mathrm{V} \cap \mathrm{Ker}(u) = \overline{v}(\{0\})\) 是紧的（TG, I, p. 75, 定理 1）；于是向量空间 \(\mathrm{Ker}(u)\) 是局部紧的，从而有限维（EVT, I, p. 15, 定理 3）。设 \(E_{1}\) 是 \(\mathrm{Ker}(u)\) 在 E 中的一个拓扑补（III, p. 55 的命题 1）；令 \(V_{1} = E_{1} \cap V\) 。集 \(V_{1}\) 在 V 中是闭的。映射 \(u|V_{1}\) 是真映射（TG, I, p. 74, 推论 1）且是单射，故是从 \(V_{1}\) 到 F 的闭子集上的一个同胚（TG, I, p. 72, 命题 2）。由引理 1，u 在 \(E_{1}\) 上的限制是从 \(E_{1}\) 到 F 的闭子空间上的一个同胚，故 u 是具有闭像的严格态射。
\end{enumerate}

